\subsection{Neighboring operators and commuting bond products}

\begin{definition}[Neighboring operators]%
    \label{def:mpdo_sector_eta}
    \lean{MPOTensor.sectorEta}
    \leanok
    \uses{def:mpdo_sector_tensors}
    Fix outer virtual indices $\alpha_1,\beta_3$ with the chosen Hayashi
    decomposition, and let $l_k,r_k$ be the sector tensors of
    Definition~\ref{def:mpdo_sector_tensors}. For each pair of sectors
    $(k,h)$, the neighboring operator contracts the right sector tensor of
    one site with the left sector tensor of the following site over the
    virtual bond they share:
    \[
        \eta_{k,h}=\sum_{\gamma}(r_k)_{\gamma}\otimes(l_h)_{\gamma}.
    \]
    It acts on the neighboring bond space $B_k^{R}\otimes B_h^{L}$
    (Figure~\ref{fig:mpdo_eta_sector_decomposition}). This is the operator
    displayed in
    \cite[Appendix~C.2, lines~1441--1445]{Cirac2016MPDO_arXiv}.
\end{definition}

\begin{theorem}[Neighboring bond contraction]
    \label{thm:mpdo_neighboring_bond_contraction}
    \lean{MPOTensor.physicalSlice_neighboring_contraction}
    \leanok
    \uses{def:mpdo_sector_eta}
    In the setting of Theorem~\ref{thm:mpdo_sector_factorization}, for all
    sectors $k,h$ and virtual indices $\beta_1,\alpha_3$, contracting the
    shared virtual bond of two neighboring factorized slices leaves the
    neighboring operator between the outer sector tensors:
    \[
        \begin{aligned}
        &\sum_{\gamma}
        \left[U_B\kappa_{\beta_1,\gamma}U_B^\dagger
        \right]_{(k,l_1,r_1),(k,l_1',r_1')} \\
        &\qquad {}\cdot
        \left[U_B\kappa_{\gamma,\alpha_3}U_B^\dagger
        \right]_{(h,l_2,r_2),(h,l_2',r_2')} \\
        &\quad =
        \left[(l_k)_{\beta_1}\right]_{l_1,l_1'}
        \left[\eta_{k,h}\right]_{(r_1,l_2),(r_1',l_2')} \\
        &\qquad {}\cdot
        \left[(r_h)_{\alpha_3}\right]_{r_2,r_2'}.
        \end{aligned}
    \]
    This is the single-bond step of the identity assembling
    $\bigotimes_{n}\eta_{k_n,k_{n+1}}$ from the factorized chain in
    \cite[Appendix~C.2, lines~1446--1449]{Cirac2016MPDO_arXiv}.
\end{theorem}

\begin{proof}
    \leanok
    \uses{thm:mpdo_sector_factorization}
    Substituting the sector factorization
    (Theorem~\ref{thm:mpdo_sector_factorization}) into each slice, the
    outer sector tensors decouple from the bond sum, and the inner sum over
    the shared bond $\gamma$ reads
    \[
        \sum_{\gamma}\bigl[(r_k)_{\gamma}\bigr]_{r_1,r_1'}\,
        \bigl[(l_h)_{\gamma}\bigr]_{l_2,l_2'}
        =\bigl[\eta_{k,h}\bigr]_{(r_1,l_2),(r_1',l_2')}
    \]
    by definition of the neighboring operator.
\end{proof}

\begin{definition}[Source-minimal physical-sector factorization]
    \label{def:mpdo_physical_sector_factorization}
    \lean{MPOTensor.PhysicalSectorFactorization,
      MPOTensor.PhysicalSectorFactorization.transformedPhysicalSlice,
      MPOTensor.PhysicalSectorFactorization.transformedPhysicalSlice_eq,
      MPOTensor.PhysicalSectorFactorization.sectorFinEquiv,
      MPOTensor.PhysicalSectorFactorization.sectorCoordinateTensor,
      MPOTensor.PhysicalSectorFactorization.neighborIndex_nonempty}
    \lean{MPOTensor.PhysicalSectorFactorization.sectorCoordinateTensor_apply,
      MPOTensor.PhysicalSectorFactorization.sectorCoordinateTensor_apply_same,
      MPOTensor.PhysicalSectorFactorization.sectorCoordinateTensor_apply_ne}
    \lean{MPOTensor.PhysicalSectorFactorization.transformedPhysicalSlice_apply_same,
      MPOTensor.PhysicalSectorFactorization.transformedPhysicalSlice_apply_ne}
    \leanok
    \uses{def:mpo_tensor}
    A physical-sector factorization of an MPO tensor ${\cal K}$ consists of a
    finite decomposition
    \[
      \mathbb C^d\simeq\bigoplus_k(B_k^L\otimes B_k^R),
    \]
    whose left and right factors are nonzero,
    an isometry $U$ on the physical space, and matrix families
    $(l_k)_\beta$ and $(r_k)_\alpha$ such that every physical slice satisfies
    \[
      U\kappa_{\beta,\alpha}U^\dagger
      =\bigoplus_k(l_k)_\beta\otimes(r_k)_\alpha.
    \]
    No positivity, trace factorization, or quantum-Markov decomposition is
    included in this datum. This is the factorization displayed in
    \cite[Appendix~C.2, lines~1381--1388]{Cirac2016MPDO_arXiv}.
\end{definition}

\begin{definition}[Algebraic neighboring contraction]
    \label{def:mpdo_minimal_neighboring_operator}
    \lean{MPOTensor.PhysicalSectorFactorization.neighboringOperator,
      MPOTensor.PhysicalSectorFactorization.neighboringOperator_apply}
    \leanok
    \uses{def:mpdo_physical_sector_factorization}
    For sectors $k,h$, contract the common virtual index of the right and
    left sector tensors:
    \[
      \eta_{k,h}
      =\sum_a(r_k)_a\otimes(l_h)_a.
    \]
    This is an operator on $B_k^R\otimes B_h^L$. Positivity is not part of
    this definition; it is a separate hypothesis in Proposition~C.7.
\end{definition}

\begin{definition}[One-site matrix spaces of sector sets]%
    \label{def:mpdo_one_site_sector_matrix_space}
    \lean{MPOTensor.PhysicalSectorFactorization.oneSiteSectorMatrix,
      MPOTensor.PhysicalSectorFactorization.oneSiteSectorSpace}
    \leanok
    \uses{def:mpdo_physical_sector_factorization}
    For a sector $k$ and physical matrix indices $x,y$ in its summand, let
    $A_{k;x,y}\in\MN{D}$ be the corresponding entry of the
    sector-coordinate tensor. For a set $S$ of sectors, define
    \[
        {\cal A}_S=\spn\{A_{k;x,y}:k\in S\}.
    \]
    These are the matrix spaces associated with the physical projections in
    \cite[Appendix~C.2, lines~1463--1467]{Cirac2016MPDO_arXiv}.
\end{definition}

\begin{theorem}[Directed-cut annihilation of one-site matrix spaces]%
    \label{thm:mpdo_directed_cut_one_site_spaces}
    \lean{MPOTensor.PhysicalSectorFactorization.%
      oneSiteSectorMatrix_mul_eq_zero_of_neighboringOperator_eq_zero}
    \lean{MPOTensor.PhysicalSectorFactorization.oneSiteSectorSpace_mul_eq_zero}
    \leanok
    \uses{def:mpdo_minimal_neighboring_operator,
      def:mpdo_one_site_sector_matrix_space}
    Let $S$ and $C$ be sets of sectors. If
    \[
        \eta_{k,h}=0\qquad(k\in S,\ h\in C),
    \]
    then
    \[
        XY=0\qquad(X\in{\cal A}_S,\ Y\in{\cal A}_C).
    \]
    No conclusion in the reverse product order is asserted. This is the
    directed local-orthogonality implication used in
    \cite[Appendix~C.2, lines~1463--1470]{Cirac2016MPDO_arXiv}.
\end{theorem}

\begin{proof}
    \leanok
    Expand two generating matrices in sector coordinates. Their product has
    entries
    \[
      \sum_{\gamma}
        (l_k)_\beta (r_k)_\gamma
        (l_h)_\gamma (r_h)_\alpha.
    \]
    The middle contraction is a matrix entry of $\eta_{k,h}$ and therefore
    vanishes. Bilinearity extends the conclusion from the generators to the
    two linear spans.
\end{proof}

\begin{definition}[Vertex rephasing of the sector tensors]%
    \label{def:mpdo_physical_sector_rephasing}
    \lean{MPOTensor.PhysicalSectorFactorization.rephase}
    \leanok
    \uses{def:mpdo_physical_sector_factorization,
      def:mpdo_minimal_neighboring_operator}
    Let $z_k\in\mathbb C$ satisfy $|z_k|=1$ for every sector. Replacing
    \[
        (l_k)_a\longmapsto z_k(l_k)_a,
        \qquad
        (r_k)_a\longmapsto z_k^{-1}(r_k)_a
    \]
    leaves every physical-slice factorization unchanged.
\end{definition}

\begin{theorem}[Rephasing identities for the sector data]%
    \label{thm:mpdo_physical_sector_rephasing_invariance}
    \lean{MPOTensor.PhysicalSectorFactorization.rephase_neighboringOperator}
    \lean{MPOTensor.PhysicalSectorFactorization.rephase_cyclicNeighboringProduct}
    \lean{MPOTensor.PhysicalSectorFactorization.rephase_sectorVirtualMatrix}
    \lean{MPOTensor.PhysicalSectorFactorization.%
      rephase_neighboringOperator_ne_zero_iff}
    \leanok
    \uses{def:mpdo_physical_sector_rephasing,
      def:mpdo_minimal_neighboring_operator}
    Vertex rephasing transforms the neighboring operators by
    \[
        \eta_{k,h}\longmapsto z_k^{-1}z_h\,\eta_{k,h}.
    \]
    It leaves every sector virtual matrix unchanged, preserves the nonzero
    neighboring support, and leaves every complete cyclic product of
    neighboring operators unchanged. These are the rephasing identities used
    in \cite[Appendix~C.2, lines~1434--1450]{Cirac2016MPDO_arXiv}.
\end{theorem}

\begin{definition}[Vanishing right tensors on inactive sectors]
    \label{def:mpdo_physical_sector_inactive_right_zero}
    \lean{MPOTensor.PhysicalSectorFactorization.zeroRightTensorOnInactive}
    \lean{MPOTensor.PhysicalSectorFactorization.zeroRightTensorOnInactive_neighboringOperator}
    \leanok
    \uses{def:mpdo_physical_sector_factorization,
      def:mpdo_minimal_neighboring_operator}
    Let $S$ be a set of active sectors in a physical-sector factorization,
    and suppose that $(l_k)_\beta=0$ for every $k\notin S$ and every $\beta$.
    Retain every physical summand and define
    \[
        \widehat l_k=l_k,
        \qquad
        \widehat r_k=
        \begin{cases}
          r_k,&k\in S,\\
          0,&k\notin S.
        \end{cases}
    \]
    Then $\widehat l_k\otimes\widehat r_k=l_k\otimes r_k$ for every $k$,
    so the physical-slice factorization is unchanged. The neighboring
    operators become
    \[
        \widehat\eta_{k,h}=
        \begin{cases}
          \eta_{k,h},&k\in S,\\
          0,&k\notin S.
        \end{cases}
    \]
    This is the reparameterization freedom in the factorization and
    neighboring contraction of
    \cite[Appendix~C.2, lines~1434--1445]{Cirac2016MPDO_arXiv}.
\end{definition}

\begin{definition}[Inverse-map physical-sector factorization]
    \label{def:mpdo_inverse_map_physical_sector_factorization}
    \lean{MPOTensor.inverseMapPhysicalSectorFactorization}
    \leanok
    \uses{thm:mpdo_sector_factorization,
      def:mpdo_physical_sector_factorization,
      lem:mpdo_trace_one_nonempty}
    Let ${\cal K}$ be injective, let $\rho$ be its three-site closure against
    a virtual tail $R$, and choose a Hayashi decomposition of $\rho$. Fix
    outer indices such that $R_{\beta_3,\alpha_1}\ne0$. The inverse-map
    tensors $l_k,r_k$, together with the Hayashi physical decomposition and
    unitary, determine a physical-sector factorization of ${\cal K}$.

    The left and right factor spaces are nonzero because the corresponding
    Hayashi density matrices have trace one; their nonzero dimensions are
    conclusions rather than hypotheses. No positivity of the neighboring
    operators is asserted. This gives the factorization in
    \cite[Appendix~C.2, equations \texttt{AppUkU=rl} and \texttt{formK},
    lines~1383--1387 and~1434--1439]{Cirac2016MPDO_arXiv}.
\end{definition}

\begin{theorem}[Identification of the neighboring operators]
    \label{thm:mpdo_inverse_map_factorization_neighboring_operator}
    \lean{MPOTensor.inverseMapPhysicalSectorFactorization_neighboringOperator_eq_sectorEta}
    \leanok
    \uses{def:mpdo_inverse_map_physical_sector_factorization,
      def:mpdo_sector_eta,
      def:mpdo_minimal_neighboring_operator}
    For all sectors $k,h$, the neighboring contraction of the resulting
    physical-sector factorization is the inverse-map operator
    \[
      \eta_{k,h}=\sum_\gamma (r_k)_\gamma\otimes(l_h)_\gamma.
    \]
    This is equation \texttt{etarl} in
    \cite[Appendix~C.2, lines~1441--1445]{Cirac2016MPDO_arXiv}.
\end{theorem}

\begin{definition}[Zero-weight inverse-map reparameterization]
    \label{def:mpdo_zero_weight_reparameterized_inverse_map_factorization}
    \lean{MPOTensor.zeroWeightReparameterizedInverseMapPhysicalSectorFactorization}
    \leanok
    \uses{def:mpdo_inverse_map_physical_sector_factorization,
      def:mpdo_physical_sector_inactive_right_zero,
      lem:mpdo_sector_tensor_left_zero_of_weight_zero}
    In the inverse-map factorization, retain every physical sector and its
    factor spaces, and set
    \[
        \widehat l_k=l_k,
        \qquad
        \widehat r_k=
        \begin{cases}
          r_k,&p_k\ne0,\\
          0,&p_k=0.
        \end{cases}
    \]
    The physical-slice factorization is unchanged. This treats the
    zero-weight freedom left by the Hayashi decomposition used in
    \cite[Appendix~C.2, lines~1415--1445]{Cirac2016MPDO_arXiv}.
\end{definition}

\begin{theorem}[Neighboring operators after zero-weight reparameterization]
    \label{thm:mpdo_zero_weight_reparameterized_neighboring_operator}
    \lean{MPOTensor.zeroWeightReparameterizedInverseMapPhysicalSectorFactorization_neighboringOperator}
    \lean{MPOTensor.probability_ne_zero_of_reparameterized_neighboringOperator_ne_zero}
    \leanok
    \uses{def:mpdo_zero_weight_reparameterized_inverse_map_factorization,
      thm:mpdo_inverse_map_factorization_neighboring_operator}
    The neighboring operators of the reparameterized factorization are
    \[
        \widehat\eta_{k,h}=
        \begin{cases}
          \eta_{k,h},&p_k\ne0,\\
          0,&p_k=0.
        \end{cases}
    \]
    Thus a zero-weight sector has no outgoing support edge. Since
    $p_h=0$ also gives $l_h=0$, it has no incoming support edge either.
\end{theorem}

\begin{proof}
    \leanok
    The neighboring contraction uses $r_k$ at its source and $l_h$ at its
    target. The first assertion follows from the definition of
    $\widehat r_k$, and the last assertion follows from
    Lemma~\ref{lem:mpdo_sector_tensor_left_zero_of_weight_zero}.
\end{proof}

\begin{corollary}[SAL gives a raw physical-sector factorization]
    \label{cor:mpdo_sal_raw_physical_sector_factorization}
    \lean{MPOTensor.nonempty_physicalSectorFactorization_of_isSAL}
    \leanok
    \uses{thm:mpdo_four_site_sal_eta_structure,
      thm:mpdo_normalized_four_site_tail_entry_ne_zero,
      def:mpdo_inverse_map_physical_sector_factorization}
    Every injective MPO tensor satisfying SAL admits a physical-sector
    factorization. This is only the raw algebraic factorization from Lemma C.4
    and equations \texttt{formK} and \texttt{etarl}; coherent positivity of
    the neighboring operators, recurrence, and primitivity are not part of
    this conclusion.
\end{corollary}

\begin{definition}[Finite-chain coordinates of a physical-sector factorization]
    \label{def:mpdo_physical_sector_chain_coordinates}
    \lean{MPOTensor.PhysicalSectorFactorization.SectorChainFiber,
      MPOTensor.PhysicalSectorFactorization.SectorChainIndex,
      MPOTensor.PhysicalSectorFactorization.sectorChainEquiv,
      MPOTensor.PhysicalSectorFactorization.cyclicNeighboringProduct}
    \leanok
    \uses{def:mpdo_physical_sector_factorization,
      def:mpdo_minimal_neighboring_operator}
    Let $N\geq1$. For a cyclic sector configuration
    $k=(k_0,\ldots,k_{N-1})$, write
    \[
      I_k=\prod_{n=0}^{N-1}(B_{k_n}^L\times B_{k_n}^R).
    \]
    Applying the one-site sector decomposition at every site gives a
    bijection
    \[
      \Phi_N:\{0,\ldots,d-1\}^N
        \longrightarrow\coprod_{k_0,\ldots,k_{N-1}}I_k.
    \]
    On $I_k$, define the cyclic neighboring product by
    \[
      \Omega_k(x,y)=\prod_{n=0}^{N-1}
        \eta_{k_n,k_{n+1}}
        \bigl((x_n^R,x_{n+1}^L),(y_n^R,y_{n+1}^L)\bigr),
    \]
    where site labels are read modulo $N$.
\end{definition}

\begin{theorem}[Finite-chain decomposition of a physical-sector factorization]
    \label{thm:mpdo_physical_sector_chain_decomposition}
    \lean{MPOTensor.PhysicalSectorFactorization.mpo_submatrix_sector_eq_cyclicNeighboringProduct,
      MPOTensor.PhysicalSectorFactorization.mpo_reindex_sectorChainEquiv_eq_blockDiagonal}
    \leanok
    \uses{def:mpdo_physical_sector_chain_coordinates}
    Suppose that the physical-sector factorization is given. For every
    $N\geq1$, the physically transformed MPO is block diagonal in the cyclic
    sector configurations, and
    \[
      \operatorname{reindex}_{\Phi_N}
        \bigl(M_N(\widetilde{\cal K})\bigr)
      =\bigoplus_{k_0,\ldots,k_{N-1}}\Omega_k.
    \]
    This is the all-length algebraic decomposition in
    \cite[Appendix~C.2, lines~1435--1450]{Cirac2016MPDO_arXiv}, conditional
    here on the displayed physical-sector factorization. It does not assert
    positivity of the individual neighboring operators.
\end{theorem}

\begin{proof}
    \leanok
    Fix a sector configuration $k$ and entries $x,y\in I_k$. Expanding the
    closed horizontal contraction gives
    \[
      \sum_g\prod_n
        (l_{k_n})_{g_n}(x_n^L,y_n^L)
        (r_{k_n})_{g_{n+1}}(x_n^R,y_n^R).
    \]
    Cyclically shifting the summation variables and collecting the factors
    with the same horizontal index turns this expression into
    $\Omega_k(x,y)$. If the row and column sector configurations differ,
    one transformed physical slice lies between distinct direct summands and
    vanishes. Thus all off-diagonal sector blocks vanish.
\end{proof}

\begin{corollary}[SAL gives positive cyclic sector products]
    \label{cor:mpdo_sal_positive_cyclic_sector_products}
    \lean{MPOTensor.exists_physicalSectorFactorization_with_cyclicNeighboringProduct_posSemidef_of_isSAL}
    \leanok
    \uses{cor:mpdo_sal_raw_physical_sector_factorization,
      thm:mpdo_conjugated_chain_congruence,
      thm:mpdo_physical_sector_chain_decomposition}
    Every injective MPO tensor satisfying SAL admits a physical-sector
    factorization such that, for every $N\geq1$ and every cyclic sector
    configuration $k=(k_0,\ldots,k_{N-1})$,
    \[
      \bigotimes_{n=0}^{N-1}\eta_{k_n,k_{n+1}}\geq0,
      \qquad k_N=k_0.
    \]
    This is the projected-chain inequality in
    \cite[Appendix~C.2, Lemma~C.4, lines~1446--1450]{Cirac2016MPDO_arXiv}.
    It does not assert positivity of the individual neighboring operators.
\end{corollary}

\begin{proof}
    \leanok
    \uses{cor:mpdo_sal_raw_physical_sector_factorization,
      thm:mpdo_conjugated_chain_congruence,
      thm:mpdo_physical_sector_chain_decomposition}
    Choose the raw physical-sector factorization supplied by SAL. The physical
    basis change sends the positive $N$-site MPO to a positive congruence.
    Each cyclic neighboring product is a diagonal sector compression of this
    congruence and is therefore positive semidefinite.
\end{proof}

\begin{definition}[Boundary contraction against a virtual matrix]
    \label{def:mpdo_sector_boundary_operator}
    \lean{MPOTensor.PhysicalSectorFactorization.boundaryOperator,
      MPOTensor.PhysicalSectorFactorization.boundaryOperator_apply,
      MPOTensor.PhysicalSectorFactorization.boundaryOperator_zero,
      MPOTensor.PhysicalSectorFactorization.boundaryOperator_add,
      MPOTensor.PhysicalSectorFactorization.boundaryOperator_smul}
    \leanok
    \uses{def:mpdo_physical_sector_factorization}
    For an arbitrary virtual matrix $X$, define the outer boundary operator
    on $B_k^L\otimes B_h^R$ by
    \[
      B_{k,h}(X)
      =\sum_{a,b}X_{b,a}\,(l_k)_a\otimes(r_h)_b.
    \]
    The map $X\mapsto B_{k,h}(X)$ is complex-linear. No positivity of $X$ or
    $B_{k,h}(X)$ is assumed.
\end{definition}

\begin{definition}[Regrouped two-site sector closure]
    \label{def:mpdo_two_site_sector_closure}
    \lean{MPOTensor.PhysicalSectorFactorization.twoSiteRegroupEquiv,
      MPOTensor.PhysicalSectorFactorization.twoSiteRegroupEquiv_symm_apply,
      MPOTensor.PhysicalSectorFactorization.twoSiteSectorEmbedding,
      MPOTensor.PhysicalSectorFactorization.twoSiteSectorClosure}
    \leanok
    \uses{def:phys_close, def:mpdo_physical_sector_factorization}
    First pass to the physical basis selected by $U$, so that each slice is
    written in the sector coordinates of
    Definition~\ref{def:mpdo_physical_sector_factorization}. Form the
    two-site closure of this sector-coordinate tensor, restrict it to the
    physical sectors $(k,h)$, and reindex its physical factors by
    \[
      ((L_k,R_k),(L_h,R_h))
      \longmapsto ((L_k,R_h),(R_k,L_h)).
    \]
    Denote the resulting matrix by ${\cal K}_{2;k,h}(X)$.
\end{definition}

\begin{theorem}[Two-site closure factorization]
    \label{thm:mpdo_two_site_sector_closure_factorization}
    \lean{MPOTensor.PhysicalSectorFactorization.twoSiteSectorClosure_eq}
    \leanok
    \uses{def:mpdo_two_site_sector_closure,
      def:mpdo_minimal_neighboring_operator,
      def:mpdo_sector_boundary_operator}
    For every virtual matrix $X$ and every sector pair $(k,h)$,
    \[
      {\cal K}_{2;k,h}(X)
      =B_{k,h}(X)\otimes\eta_{k,h}.
    \]
\end{theorem}
\begin{proof}
    \leanok
    \uses{def:mpdo_two_site_sector_closure,
      def:mpdo_minimal_neighboring_operator,
      def:mpdo_sector_boundary_operator,
      def:phys_close,
      def:mpdo_physical_sector_factorization}
    At the regrouped matrix entry indexed by
    $((\lambda_k,\rho_h),(\rho_k,\lambda_h))$ and
    $((\lambda'_k,\rho'_h),(\rho'_k,\lambda'_h))$, the virtual trace expands
    as the scalar
    \[
      \sum_{a,b,c}X_{b,a}
      [(l_k)_a]_{\lambda_k,\lambda'_k}
      [(r_k)_c]_{\rho_k,\rho'_k}
      [(l_h)_c]_{\lambda_h,\lambda'_h}
      [(r_h)_b]_{\rho_h,\rho'_h}.
    \]
    Commuting scalar factors and collecting the sums over $(a,b)$ and $c$
    gives respectively $B_{k,h}(X)$ and $\eta_{k,h}$.
\end{proof}

\begin{definition}[Regrouped three-site sector closure]
    \label{def:mpdo_three_site_sector_closure}
    \lean{MPOTensor.PhysicalSectorFactorization.threeSiteRegroupEquiv,
      MPOTensor.PhysicalSectorFactorization.threeSiteRegroupEquiv_symm_apply,
      MPOTensor.PhysicalSectorFactorization.threeSiteSectorEmbedding,
      MPOTensor.PhysicalSectorFactorization.threeSiteSectorClosure}
    \leanok
    \uses{def:phys_close, def:mpdo_physical_sector_factorization}
    First pass to the physical basis selected by $U$ and form the three-site
    closure of the resulting sector-coordinate tensor. Restrict this closure
    to the sectors $(k,l,h)$ and reindex its physical factors by
    \[
      ((L_k,R_k),((L_l,R_l),(L_h,R_h)))
      \longmapsto
      ((L_k,R_h),((R_k,L_l),(R_l,L_h))).
    \]
    Denote the resulting matrix by ${\cal K}_{3;k,l,h}(X)$.
\end{definition}

\begin{theorem}[Three-site closure factorization]
    \label{thm:mpdo_three_site_sector_closure_factorization}
    \lean{MPOTensor.PhysicalSectorFactorization.threeSiteSectorClosure_eq}
    \leanok
    \uses{def:mpdo_three_site_sector_closure,
      def:mpdo_minimal_neighboring_operator,
      def:mpdo_sector_boundary_operator}
    For every virtual matrix $X$ and fixed sectors $(k,l,h)$,
    \[
      {\cal K}_{3;k,l,h}(X)
      =B_{k,h}(X)\otimes
        \bigl(\eta_{k,l}\otimes\eta_{l,h}\bigr).
    \]
\end{theorem}
\begin{proof}
    \leanok
    \uses{def:mpdo_three_site_sector_closure,
      def:mpdo_minimal_neighboring_operator,
      def:mpdo_sector_boundary_operator,
      def:phys_close,
      def:mpdo_physical_sector_factorization}
    At the regrouped matrix entry indexed by
    $((\lambda_k,\rho_h),((\rho_k,\lambda_l),(\rho_l,\lambda_h)))$ and its
    primed counterpart, the left-hand side is the scalar
    \[
      \sum_{a,b,c,e}X_{b,a}
      [(l_k)_a]_{\lambda_k,\lambda'_k}
      [(r_k)_c]_{\rho_k,\rho'_k}
      [(l_l)_c]_{\lambda_l,\lambda'_l}
      [(r_l)_e]_{\rho_l,\rho'_l}
      [(l_h)_e]_{\lambda_h,\lambda'_h}
      [(r_h)_b]_{\rho_h,\rho'_h}.
    \]
    Reordering the four finite sums separates the $(a,b)$ boundary
    contraction from the neighboring contractions indexed by $c$ and $e$,
    which gives the three factors on the right-hand side.
\end{proof}

\begin{definition}[Three-site neighboring operator]
    \label{def:mpdo_physical_sector_three_site_neighboring_operator}
    \lean{MPOTensor.PhysicalSectorFactorization.threeSiteNeighboringOperator}
    \leanok
    \uses{def:mpdo_minimal_neighboring_operator}
    For every outer-sector pair $(k,h)$, define
    \[
      \Omega_{k,h}=\bigoplus_l
        \bigl(\eta_{k,l}\otimes\eta_{l,h}\bigr).
    \]
    This is the direct sum of the two neighboring contractions appearing in
    the three-site closure factorization
    \cite[Appendix~C.2, lines~1510--1516]{Cirac2016MPDO_arXiv}.
\end{definition}

\begin{figure}[ht]
    \centering
    \begin{minipage}{0.41\textwidth}
      \centering
      \TNMPDOTwoSiteClosureFactorization
    \end{minipage}\hfill
    \begin{minipage}{0.56\textwidth}
      \centering
      \TNMPDOThreeSiteClosureFactorization
    \end{minipage}
    \caption{First each slice is conjugated by $U$ and expressed in the
    corresponding sector coordinates. After restriction to fixed sectors and
    the displayed regroupings, the arbitrary-$X$ two- and three-site closures
    separate into the boundary factor and respectively one or two neighboring
    contractions. The diagrams do not assert positivity,
    preparation, recovery, or the full conclusion of Proposition~C.7. They
    follow the contractions in
    \cite[lines~638--655 and Appendix~C.2, lines~1435--1448]
      {Cirac2016MPDO_arXiv}.}
    \label{fig:mpdo_sector_closure_factorizations}
\end{figure}

\begin{definition}[Neighboring-operator trace factorization]
    \label{def:mpdo_neighboring_trace_factorization}
    \lean{MPOTensor.PhysicalSectorFactorization.NeighboringTraceFactorization}
    \lean{MPOTensor.PhysicalSectorFactorization.NeighboringTraceFactorization.weight_nonneg,
      MPOTensor.PhysicalSectorFactorization.NeighboringTraceFactorization.neighboringOperator_eq_zero_of_mul_eq_zero,
      MPOTensor.PhysicalSectorFactorization.NeighboringTraceFactorization.neighborIndex_nonempty_of_mul_ne_zero}
    \leanok
    \uses{def:mpdo_minimal_neighboring_operator}
    Suppose that every neighboring operator is positive semidefinite and that
    there are real numbers $a_k,b_h$ such that
    \[
      \tr(\eta_{k,h})=a_kb_h,
      \qquad
      \sum_k a_kb_k=1.
    \]
    These are precisely the remaining hypotheses used in the construction of
    Proposition~C.7 after the physical-sector factorization has been fixed
    \cite[Appendix~C.2, lines~1389--1403]{Cirac2016MPDO_arXiv}.
\end{definition}

\begin{theorem}[Trace of the three-site neighboring operator]
    \label{thm:mpdo_physical_sector_three_site_neighboring_operator_trace}
    \lean{MPOTensor.PhysicalSectorFactorization.trace_threeSiteNeighboringOperator,
      MPOTensor.PhysicalSectorFactorization.NeighboringTraceFactorization.trace_threeSiteNeighboringOperator}
    \leanok
    \uses{def:mpdo_physical_sector_three_site_neighboring_operator,
      def:mpdo_neighboring_trace_factorization}
    Assume the neighboring-operator trace factorization of
    Definition~\ref{def:mpdo_neighboring_trace_factorization}. Then, for every
    outer-sector pair $(k,h)$,
    \[
      \tr(\Omega_{k,h})
      =\sum_l\tr(\eta_{k,l})\tr(\eta_{l,h})
      =a_kb_h.
    \]
    This is the normalization used in the coarse-graining channel of
    \cite[Appendix~C.2]{Cirac2016MPDO_arXiv}, lines~1547--1555.
\end{theorem}

\begin{proof}
    \leanok
    \uses{def:mpdo_neighboring_trace_factorization}
    The trace of a direct sum is the sum of the traces of its blocks, and the
    trace of a tensor product is the product of the traces. Hence
    \[
      \tr(\Omega_{k,h})
      =\sum_l(a_kb_l)(a_lb_h)
      =a_kb_h\sum_l a_lb_l
      =a_kb_h.
    \]
\end{proof}

\begin{theorem}[Partial traces of the sector closures]
    \label{thm:mpdo_sector_closure_partial_traces}
    \lean{MPOTensor.PhysicalSectorFactorization.partialTraceRight_twoSiteSectorClosure}
    \lean{MPOTensor.PhysicalSectorFactorization.NeighboringTraceFactorization.partialTraceRight_twoSiteSectorClosure}
    \lean{MPOTensor.PhysicalSectorFactorization.partialTraceRight_threeSiteSectorClosure}
    \lean{MPOTensor.PhysicalSectorFactorization.NeighboringTraceFactorization.sum_threeSite_trace_coefficients}
    \leanok
    \uses{thm:mpdo_two_site_sector_closure_factorization,
      thm:mpdo_three_site_sector_closure_factorization,
      def:mpdo_neighboring_trace_factorization}
    For every virtual matrix $X$,
    \[
      \tr_{R_k\otimes L_h}\!\left({\cal K}_{2;k,h}(X)\right)
      =\tr(\eta_{k,h})B_{k,h}(X)
      =a_kb_h B_{k,h}(X).
    \]
    Similarly,
    \[
      \tr_{(R_k\otimes L_l)\otimes(R_l\otimes L_h)}
        \!\left({\cal K}_{3;k,l,h}(X)\right)
      =\tr(\eta_{k,l})\tr(\eta_{l,h})B_{k,h}(X),
    \]
    and summing the scalar coefficient over $l$ gives $a_kb_h$.
\end{theorem}

\begin{proof}
    \leanok
    \uses{thm:mpdo_two_site_sector_closure_factorization,
      thm:mpdo_three_site_sector_closure_factorization}
    The partial trace of $A\otimes B$ is $\tr(B)A$. The two-site identity
    follows immediately. For three sites, the coefficient is
    \[
      \sum_l (a_kb_l)(a_lb_h)
      =a_k\left(\sum_l a_lb_l\right)b_h
      =a_kb_h.
    \]
\end{proof}

\begin{definition}[Global two- and three-site sector coordinates]
    \label{def:mpdo_global_sector_closure_coordinates}
    \lean{MPOTensor.PhysicalSectorFactorization.twoSiteGlobalRegroupEquiv,
      MPOTensor.PhysicalSectorFactorization.twoSiteGlobalRegroupEquiv_symm_apply,
      MPOTensor.PhysicalSectorFactorization.threeSiteGlobalRegroupEquiv,
      MPOTensor.PhysicalSectorFactorization.threeSiteGlobalRegroupEquiv_symm_apply}
    \leanok
    \uses{def:mpdo_two_site_sector_closure,
      def:mpdo_three_site_sector_closure}
    The two- and three-site physical spaces are reindexed by their canonical
    direct sums over sector pairs and triples, with each summand further
    regrouped into its boundary and neighboring factors. Write
    $\mathfrak G_2$ and $\mathfrak G_3$ for these reindexings of matrices.
\end{definition}

\begin{theorem}[Global direct-sum form of the sector closures]
    \label{thm:mpdo_global_sector_closure_decomposition}
    \lean{MPOTensor.PhysicalSectorFactorization.twoSiteGlobalClosure_eq,
      MPOTensor.PhysicalSectorFactorization.threeSiteGlobalClosure_eq}
    \leanok
    \uses{thm:mpdo_two_site_sector_closure_factorization,
      thm:mpdo_three_site_sector_closure_factorization,
      def:mpdo_physical_sector_factorization,
      def:mpdo_global_sector_closure_coordinates}
    After the canonical regrouping of all physical-sector indices, the full
    closures are
    \[
      \mathfrak G_2({\cal K}_2(X))
      =\bigoplus_{k,h}B_{k,h}(X)\otimes\eta_{k,h}
    \]
    and
    \[
      \mathfrak G_3({\cal K}_3(X))
      =\bigoplus_{k,l,h}B_{k,h}(X)\otimes
        \bigl(\eta_{k,l}\otimes\eta_{l,h}\bigr).
    \]
    Every matrix entry between distinct sector tuples vanishes.
\end{theorem}

\begin{proof}
    \leanok
    \uses{thm:mpdo_two_site_sector_closure_factorization,
      thm:mpdo_three_site_sector_closure_factorization,
      def:mpdo_physical_sector_factorization,
      def:mpdo_global_sector_closure_coordinates}
    On a diagonal sector tuple these are the fixed-sector factorizations.
    For distinct sector labels $k\ne p$, the block structure of the
    sector-coordinate tensor gives
    \[
      [U\kappa_{\beta,\alpha}U^\dagger]_{k,p}=0.
    \]
    Consequently, if $(k,h)\ne(p,q)$ or
    $(k,l,h)\ne(p,m,q)$, respectively, then
    \[
      [\mathfrak G_2({\cal K}_2(X))]_{(k,h),(p,q)}=0,
      \qquad
      [\mathfrak G_3({\cal K}_3(X))]_{(k,l,h),(p,m,q)}=0.
    \]
\end{proof}

\begin{definition}[Two-site sector bond]
    \label{def:mpdo_physical_sector_two_site_bond}
    \lean{MPOTensor.PhysicalSectorFactorization.sectorCoordinateBond}
    \leanok
    \uses{def:mpdo_minimal_neighboring_operator,
      def:mpdo_global_sector_closure_coordinates}
    In the two-site sector coordinates, define
    \[
      B_{\mathrm{sector}}
      =\bigoplus_{k,h}
        \left(I_{B_k^L\otimes B_h^R}\otimes\eta_{k,h}\right).
    \]
    Thus the outer factors carry the identity, while the neighboring factors
    carry the operator $\eta_{k,h}$. This is the local bond in
    \cite[Appendix~C.2, Proposition~C.8, lines~1581--1593]
      {Cirac2016MPDO_arXiv}.
\end{definition}

\begin{theorem}[Positivity of the two-site sector bond]
    \label{thm:mpdo_physical_sector_two_site_bond_pos}
    \lean{MPOTensor.PhysicalSectorFactorization.sectorCoordinateBond_pos}
    \leanok
    \uses{def:mpdo_physical_sector_two_site_bond,
      thm:matrix_pos_semidef_block_diagonal}
    Suppose that $\eta_{k,h}$ is positive semidefinite for every pair of
    sectors. Then $B_{\mathrm{sector}}$ is positive semidefinite. The positivity
    premise is the conclusion at
    \cite[Appendix~C.2, lines~1446--1450]{Cirac2016MPDO_arXiv}.
\end{theorem}

\begin{proof}
    \leanok
    \uses{thm:matrix_pos_semidef_block_diagonal}
    Each summand
    $I_{B_k^L\otimes B_h^R}\otimes\eta_{k,h}$ is positive semidefinite.
    Positivity is preserved by finite block diagonals and by the reindexing
    from the regrouped direct sum to the two-site sector coordinates.
\end{proof}

\begin{definition}[Two-site bond in physical coordinates]
    \label{def:mpdo_physical_two_site_bond}
    \lean{MPOTensor.PhysicalSectorFactorization.physicalPairBond,
      MPOTensor.PhysicalSectorFactorization.physicalBond}
    \leanok
    \uses{def:mpdo_physical_sector_two_site_bond,
      def:mpdo_physical_sector_factorization}
    Let $V_2$ be the two-site tensor product of the project's coordinate map,
    so that $V_2XV_2^\dagger$ expresses a physical operator $X$ in sector
    coordinates. Thus $V_2=U_{\mathrm{paper}}^\dagger$ in the convention of
    \cite[Appendix~C.2, lines~1581--1583]{Cirac2016MPDO_arXiv}. Define
    \[
      B_{\mathrm{physical}}=V_2^\dagger B_{\mathrm{sector}}V_2,
    \]
    and identify pairs of physical indices with functions
    $\{0,1\}\to\{0,\ldots,d-1\}$. This is the physical two-site bond in
    \cite[Appendix~C.2, Proposition~C.8, lines~1581--1593]
      {Cirac2016MPDO_arXiv}.
\end{definition}

\begin{theorem}[Positivity of the physical two-site bond]
    \label{thm:mpdo_physical_two_site_bond_pos}
    \lean{MPOTensor.PhysicalSectorFactorization.physicalPairBond_pos,
      MPOTensor.PhysicalSectorFactorization.physicalBond_pos}
    \leanok
    \uses{def:mpdo_physical_two_site_bond,
      thm:mpdo_physical_sector_two_site_bond_pos}
    Suppose that $\eta_{k,h}$ is positive semidefinite for every pair of
    sectors. Then $B_{\mathrm{physical}}$ is positive semidefinite, both with
    pair indices and with two-site configuration indices.
\end{theorem}

\begin{proof}
    \leanok
    \uses{thm:mpdo_physical_sector_two_site_bond_pos}
    Apply positivity of $B_{\mathrm{sector}}$ to the congruence
    $V_2^\dagger B_{\mathrm{sector}}V_2$. Reindexing the resulting matrix by
    two-site configurations preserves positivity.
\end{proof}

\begin{definition}[Three-site lifts of a two-site operator]
    \label{def:mpdo_three_site_pair_lifts}
    \lean{MPOTensor.PhysicalSectorFactorization.leftPairMatrix,
      MPOTensor.PhysicalSectorFactorization.rightPairMatrix}
    \leanok
    Let $B$ be an operator on two copies of a finite-dimensional space $H$.
    On the right-associated space $H\otimes(H\otimes H)$, define
    \[
      B_{01}=\alpha(B\otimes I_H)\alpha^{-1},
      \qquad
      B_{12}=I_H\otimes B,
    \]
    where $\alpha:(H\otimes H)\otimes H\to H\otimes(H\otimes H)$ is the
    canonical associator.
\end{definition}

\begin{theorem}[Three-site translation identities]
    \label{thm:mpdo_three_site_pair_lifts}
    \lean{MPOTensor.PhysicalSectorFactorization.reindex_embedLocalOperator_zero,
      MPOTensor.PhysicalSectorFactorization.reindex_embedLocalOperator_one}
    \leanok
    \uses{def:mpdo_physical_two_site_bond,
      def:mpdo_three_site_pair_lifts}
    Under the canonical identification of three-site configurations with
    $H\otimes(H\otimes H)$, the translates of the physical bond beginning at
    sites $0$ and $1$ are respectively
    \[
      (B_{\mathrm{physical}})_{01}=B_{01},
      \qquad
      (B_{\mathrm{physical}})_{12}=B_{12}.
    \]
\end{theorem}

\begin{proof}
    \leanok
    An entry of the first translate vanishes unless the third indices agree;
    its remaining entry is that of $B_{\mathrm{physical}}$. Similarly, an
    entry of the second translate vanishes unless the first indices agree.
    These are precisely the matrix entries of $B_{01}$ and $B_{12}$.
\end{proof}

\begin{theorem}[Commutativity on a fixed sector triple]
    \label{thm:mpdo_fixed_sector_adjacent_bonds_comm}
    \lean{MPOTensor.PhysicalSectorFactorization.adjacentSectorBonds_comm}
    \leanok
    \uses{def:mpdo_minimal_neighboring_operator}
    Fix sectors $(k,l,h)$. On the regrouped sector space, set
    \begin{align*}
      C_{01}^{k,l,h}
        &=I_{B_k^L\otimes B_h^R}\otimes
          \left(\eta_{k,l}\otimes I_{B_l^R\otimes B_h^L}\right),\\
      C_{12}^{k,l,h}
        &=I_{B_k^L\otimes B_h^R}\otimes
          \left(I_{B_k^R\otimes B_l^L}\otimes\eta_{l,h}\right).
    \end{align*}
    Then
    \[
      C_{01}^{k,l,h}C_{12}^{k,l,h}
        =C_{12}^{k,l,h}C_{01}^{k,l,h}.
    \]
    This is the fixed-sector calculation in
    \cite[Appendix~C.2, Proposition~C.8, lines~1589--1593]
      {Cirac2016MPDO_arXiv}.
\end{theorem}

\begin{figure}[ht]
    \centering
    \TNMPDOFixedSectorAdjacentCommutativity
    \caption{On the sector triple $(k,l,h)$, the operator
    $\eta_{k,l}$ acts only on $B_k^R\otimes B_l^L$, whereas $\eta_{l,h}$ acts
    only on $B_l^R\otimes B_h^L$. The outer factors $B_k^L$ and $B_h^R$ are
    unchanged.
    Thus exchanging the vertical order of the two operators leaves the
    contraction unchanged, which is the displayed fixed-sector commutativity.}
    \label{fig:mpdo_fixed_sector_adjacent_bonds_comm}
\end{figure}

\begin{proof}
    \leanok
    The operators $\eta_{k,l}$ and $\eta_{l,h}$ act on the distinct factors
    $B_k^R\otimes B_l^L$ and $B_l^R\otimes B_h^L$, respectively. The claim
    follows from the mixed-product identity for tensor products.
\end{proof}

\begin{theorem}[Commutativity of adjacent physical bonds]
    \label{thm:mpdo_physical_adjacent_bonds_comm}
    \lean{MPOTensor.PhysicalSectorFactorization.physicalBond_zero_one_comm}
    \leanok
    \uses{def:mpdo_physical_two_site_bond,
      def:mpdo_global_sector_closure_coordinates,
      thm:mpdo_three_site_pair_lifts,
      thm:mpdo_fixed_sector_adjacent_bonds_comm}
    The two translates of the physical bond on three sites commute:
    \[
      (B_{\mathrm{physical}})_{01}(B_{\mathrm{physical}})_{12}
        =(B_{\mathrm{physical}})_{12}(B_{\mathrm{physical}})_{01}.
    \]
    This is the adjacent-bond calculation in
    \cite[Appendix~C.2, Proposition~C.8, lines~1589--1593]
      {Cirac2016MPDO_arXiv}.
\end{theorem}

\begin{proof}
    \leanok
    Regroup the three-site sector coordinates into the direct sum over triples
    $(k,l,h)$. The two bonds become block diagonal, with blocks
    $C_{01}^{k,l,h}$ and $C_{12}^{k,l,h}$, which commute by
    Theorem~\ref{thm:mpdo_fixed_sector_adjacent_bonds_comm}. Conjugating by the
    three-fold physical coordinate unitary preserves the equality. The
    three-site translation identities then give the displayed physical
    commutator.
\end{proof}

\begin{theorem}[Agreement outside a cyclic window]
    \label{thm:mpdo_agreement_outside_cyclic_window}
    \lean{MPOTensor.agreesOutsideWindow_iff}
    \leanok
    Let $W_i^L=\{i,i+1,\ldots,i+L-1\}$ be a cyclic window of length
    $L\leq N$. Two configurations agree outside $W_i^L$ precisely when their
    values agree at every site not contained in $W_i^L$.
\end{theorem}

\begin{proof}\leanok
    Replacing the entries in $W_i^L$ does not affect any complementary site.
    Conversely, if the two configurations agree on the complement, replacing
    the window of the first by the window of the second reconstructs the
    second configuration.
\end{proof}

\begin{theorem}[Multiplicativity of cyclic local embeddings]
    \label{thm:mpdo_cyclic_local_embedding_mul}
    \lean{MPOTensor.embedLocalOperator_mul}
    \leanok
    Let $E_i^L$ embed an operator on $L$ consecutive sites into an
    $N$-site periodic chain at the cyclic window beginning at $i$. Then
    \[
        E_i^L(BC)=E_i^L(B)E_i^L(C).
    \]
\end{theorem}

\begin{proof}\leanok
    Identify a periodic configuration with its restriction to $W_i^L$ and
    its restriction to the cyclic complement. In these coordinates,
    $E_i^L(B)=B\otimes I$, so the assertion is the mixed-product identity.
\end{proof}

\begin{theorem}[Commutativity on disjoint cyclic windows]
    \label{thm:mpdo_disjoint_cyclic_local_embeddings_comm}
    \lean{MPOTensor.embedLocalOperator_commute_of_not_cyclicWindowsOverlap}
    \leanok
    \uses{lem:cyclic_window_support_offsets,
      lem:cyclic_windows_overlap_symmetry}
    Let $L\leq N$, and let $W_i^L$ and $W_j^L$ be disjoint cyclic windows in
    an $N$-site periodic chain. For arbitrary operators $B$ and $C$ on $L$ sites,
    \[
      E_i^L(B)E_j^L(C)=E_j^L(C)E_i^L(B).
    \]
\end{theorem}

\begin{proof}\leanok
    Apply both sides to a chain amplitude. The first order of application is
    a double sum over replacements on $W_i^L$ and $W_j^L$. Since the windows
    are disjoint, each replacement leaves the other restricted configuration
    unchanged, and the two replacements commute. Interchanging the two finite
    sums gives the reverse order.
\end{proof}

\begin{theorem}[Adjacent physical bonds on periodic chains]
    \label{thm:mpdo_periodic_adjacent_physical_bonds_comm}
    \lean{MPOTensor.PhysicalSectorFactorization.physicalBond_adjacent_comm}
    \leanok
    \uses{thm:mpdo_cyclic_local_embedding_mul,
      thm:mpdo_physical_adjacent_bonds_comm}
    Let $N\geq 3$. For every site $i\in\mathbb Z/N\mathbb Z$, including
    $i=N-1$, the adjacent translates of the physical bond commute:
    \[
      (B_{\mathrm{physical}})_{i,i+1}
      (B_{\mathrm{physical}})_{i+1,i+2}
      =
      (B_{\mathrm{physical}})_{i+1,i+2}
      (B_{\mathrm{physical}})_{i,i+1}.
    \]
    This is the periodic-chain form of the adjacent-bond calculation in
    \cite[Appendix~C.2, Proposition~C.8, lines~1571--1593]
      {Cirac2016MPDO_arXiv}. The case $N=2$ is treated separately below,
    because its two translated bonds have the same support with opposite
    cyclic order.
\end{theorem}

\begin{proof}
    \leanok
    Enclose the two bonds in the cyclic three-site window beginning at $i$.
    In these coordinates they are the bonds beginning at sites $0$ and $1$,
    respectively. Embedding an operator into a fixed cyclic window preserves
    products, so the assertion follows from
    Theorem~\ref{thm:mpdo_physical_adjacent_bonds_comm}.
\end{proof}

\begin{theorem}[Commutativity of the crossed two-site translates]
    \label{thm:mpdo_physical_two_site_crossed_bonds_comm}
    \lean{MPOTensor.PhysicalSectorFactorization.physicalBond_two_zero_one_comm}
    \leanok
    \uses{def:mpdo_global_sector_closure_coordinates,
      def:mpdo_physical_sector_two_site_bond,
      def:mpdo_physical_two_site_bond}
    On the periodic chain of length two, the translates beginning at sites
    zero and one read the physical sites in the opposite cyclic orders
    $(0,1)$ and $(1,0)$.  Nevertheless they commute:
    \[
      B_{01}B_{10}=B_{10}B_{01}.
    \]
    This is the length-two instance of the translated-bond commutativity in
    \cite[Appendix~C.2, Proposition~C.8]{Cirac2016MPDO_arXiv}.  The source
    does not discuss the crossed finite-size ordering separately.
\end{theorem}

\begin{proof}
    \leanok
    Fix sectors $(k,h)$ and regroup the two physical sites as
    \[
      (L_k\otimes R_h)\otimes(R_k\otimes L_h).
    \]
    The translate in the order $(0,1)$ is
    \[
      I_{L_k\otimes R_h}\otimes\eta_{k,h},
    \]
    while the translate in the order $(1,0)$ is
    \[
      \eta_{h,k}^{\mathrm{op}}\otimes I_{R_k\otimes L_h},
    \]
    where $\eta_{h,k}^{\mathrm{op}}$ denotes the same matrix after exchanging
    the factors $R_h$ and $L_k$.  These operators act on complementary
    factors and commute.  Taking the direct sum over $(k,h)$ and conjugating
    by the two-fold physical coordinate unitary proves the claim.
\end{proof}

\begin{theorem}[Disjoint physical bonds on periodic chains]
    \label{thm:mpdo_periodic_disjoint_physical_bonds_comm}
    \lean{MPOTensor.PhysicalSectorFactorization.physicalBond_disjoint_comm}
    \leanok
    \uses{def:mpdo_physical_two_site_bond,
      thm:mpdo_disjoint_cyclic_local_embeddings_comm}
    Let $N\geq 2$, and suppose that the cyclic bonds beginning at $i$ and $j$
    have disjoint two-site supports. Then their physical bond operators
    commute:
    \[
      (B_{\mathrm{physical}})_{i,i+1}
      (B_{\mathrm{physical}})_{j,j+1}
      =
      (B_{\mathrm{physical}})_{j,j+1}
      (B_{\mathrm{physical}})_{i,i+1}.
    \]
    This is the locality part of the pairwise commutation assertion in
    \cite[Appendix~C.2, Proposition~C.8, lines~1571--1593]
      {Cirac2016MPDO_arXiv}.
\end{theorem}

\begin{proof}\leanok
    This is Theorem~\ref{thm:mpdo_disjoint_cyclic_local_embeddings_comm}
    applied to the same physical bond operator in both windows. No positivity
    or additional property of the sector factorization is needed for this
    locality step.
\end{proof}

\begin{theorem}[Pairwise commutativity of the physical bonds]
    \label{thm:mpdo_periodic_pairwise_physical_bonds_comm}
    \lean{MPOTensor.PhysicalSectorFactorization.physicalBond_pairwise_comm}
    \leanok
    \uses{lem:two_site_cyclic_window_overlap_cases,
      thm:mpdo_periodic_adjacent_physical_bonds_comm,
      thm:mpdo_physical_two_site_crossed_bonds_comm,
      thm:mpdo_periodic_disjoint_physical_bonds_comm}
    Suppose that the physical-sector factorization is given. For every
    periodic chain of length $N\geq2$ and every pair of sites
    $i,j\in\mathbb Z/N\mathbb Z$, the corresponding translates of the
    physical bond commute:
    \[
      (B_{\mathrm{physical}})_{i,i+1}
      (B_{\mathrm{physical}})_{j,j+1}
      =
      (B_{\mathrm{physical}})_{j,j+1}
      (B_{\mathrm{physical}})_{i,i+1}.
    \]
    This is the pairwise commutation assertion in
    \cite[Appendix~C.2, Proposition~C.8, lines~1571--1593]
      {Cirac2016MPDO_arXiv}, conditional here on the physical-sector
    factorization constructed in the preceding results.
\end{theorem}

\begin{proof}
    \leanok
    For $N=2$, equal translates commute trivially, while the two distinct
    translates commute by
    Theorem~\ref{thm:mpdo_physical_two_site_crossed_bonds_comm}. Let $N\geq3$.
    If the two cyclic supports are disjoint, apply
    Theorem~\ref{thm:mpdo_periodic_disjoint_physical_bonds_comm}. Otherwise,
    Lemma~\ref{lem:two_site_cyclic_window_overlap_cases} shows that the two
    starting sites either coincide or are cyclic neighbors. The coincident
    case is immediate, and the two neighboring orientations follow from
    Theorem~\ref{thm:mpdo_periodic_adjacent_physical_bonds_comm} and its
    symmetric equality.
\end{proof}

\begin{theorem}[Exact product of physical-sector bonds]
    \label{thm:mpdo_exact_physical_sector_bond_product}
    \lean{MPOTensor.PhysicalSectorFactorization.mpo_eq_product_physicalBond}
    \leanok
    \uses{def:mpdo_physical_sector_factorization,
      def:mpdo_physical_two_site_bond,
      thm:mpdo_physical_sector_chain_decomposition}
    Suppose that $K$ is equipped with a physical-sector factorization $F$.
    Let $B_i$ denote the cyclic translate beginning at site $i$ of the
    physical two-site bond determined by $F$. For every $N\geq2$, the
    $N$-site periodic operator is
    \[
      \rho_N=B_0B_1\cdots B_{N-1}.
    \]
    No positivity, zero-correlation-length, injectivity, or normalization
    hypothesis is required. This is the product identity in
    \cite[Appendix~C.2, Proposition~C.8, lines~1581--1593]
      {Cirac2016MPDO_arXiv}, conditional here on the given physical-sector
    factorization.
\end{theorem}

\begin{proof}
    \leanok
    For a fixed cyclic sector sequence, contraction of the virtual index
    between the right tensor at site $i$ and the left tensor at site $i+1$
    produces $\eta_{k_i,k_{i+1}}$. Taking the virtual trace therefore gives
    the cyclic product of these neighboring operators. The
    translated sector bonds give the same ordered product. Conjugating by the
    tensor power of the one-site coordinate unitary gives the stated physical
    identity. For $N=2$, the second translate has the opposite cyclic order;
    the two factors are the ordinary and crossed bonds.
\end{proof}

\begin{theorem}[Unit scalar in the physical-sector product]
    \label{thm:mpdo_physical_sector_bond_product_unit_scalar}
    \lean{MPOTensor.PhysicalSectorFactorization.exists_positive_scalar_mpo_eq_product_physicalBond}
    \leanok
    \uses{thm:mpdo_exact_physical_sector_bond_product}
    Under the physical-sector factorization $F$ of the preceding theorem,
    for every $N\geq2$ there is a real number $c>0$ such that
    \[
      \rho_N=c\,B_0B_1\cdots B_{N-1}.
    \]
    Here the bonds are the canonical bonds determined by $F$, and one may
    take $c=1$. No positivity hypothesis on the neighboring operators
    $\eta_{k,h}$ is required.
\end{theorem}

\begin{proof}
    \leanok
    Take $c=1$ in
    Theorem~\ref{thm:mpdo_exact_physical_sector_bond_product}.
\end{proof}

\begin{theorem}[Positive translation-invariant bond datum]
    \label{thm:mpdo_positive_translation_invariant_bond_data}
    \lean{MPOTensor.PhysicalSectorFactorization.translationInvariantBondData}
    \leanok
    \uses{thm:mpdo_physical_two_site_bond_pos,
      thm:mpdo_periodic_pairwise_physical_bonds_comm}
    Suppose that every neighboring operator $\eta_{k,h}$ is positive
    semidefinite. Then the physical bond is positive semidefinite and all of
    its cyclic translates commute on every periodic chain of length at least
    two. Thus it determines a single translation-invariant positive bond
    datum. This assertion does not include the product formula for the
    finite-chain density operator.
\end{theorem}

\begin{proof}
    \leanok
    Positivity is
    Theorem~\ref{thm:mpdo_physical_two_site_bond_pos}; pairwise commutativity
    is Theorem~\ref{thm:mpdo_periodic_pairwise_physical_bonds_comm}.
\end{proof}

\begin{theorem}[$\eta$-local structure from a positive physical-sector factorization]
    \label{thm:mpdo_eta_local_structure_of_positive_physical_sector_factorization}
    \lean{MPOTensor.PhysicalSectorFactorization.etaLocalStructureData}
    \leanok
    \uses{thm:mpdo_exact_physical_sector_bond_product,
      thm:mpdo_positive_translation_invariant_bond_data,
      def:mpdo_eta_local_structure_data}
    Suppose that a physical-sector factorization of $K$ is given and that
    every neighboring operator $\eta_{k,h}$ is positive semidefinite. Then
    the physical two-site bond determines an $\eta$-local structure for $K$.
    Its translated copies commute pairwise, and for every $N\geq2$,
    \[
      \rho_N=B_0B_1\cdots B_{N-1}.
    \]
    Thus the positive normalization scalar may be chosen to be $1$. This is
    Proposition~C.8 conditional on the coherently positive physical-sector
    factorization that the source derives from SAL.
\end{theorem}

\begin{proof}
    \leanok
    Theorem~\ref{thm:mpdo_positive_translation_invariant_bond_data} gives the
    positive translation-invariant bond and its pairwise commuting
    translates. Theorem~\ref{thm:mpdo_exact_physical_sector_bond_product}
    realizes every finite-chain operator with normalization scalar $1$.
\end{proof}

\begin{corollary}[$\eta$-local structure from SAL]%
    \label{cor:mpdo_eta_local_structure_of_is_sal}
    \lean{MPOTensor.nonempty_etaLocalStructureData_of_isSAL}
    \leanok
    \uses{thm:mpdo_eta_local_structure_of_positive_physical_sector_factorization,
      thm:mpdo_eta_coherent_positive_choice}
    Every injective MPO tensor satisfying SAL admits an $\eta$-local
    structure.  Thus there is a single positive two-site bond whose cyclic
    translates commute pairwise and whose ordered product is the finite-chain
    density operator at every length $N\geq2$, with normalization scalar one.
    This is the conclusion of Proposition~C.8.
\end{corollary}

\begin{proof}
    \leanok
    Theorem~\ref{thm:mpdo_eta_coherent_positive_choice} supplies the positive
    physical-sector factorization. Apply
    Theorem~\ref{thm:mpdo_eta_local_structure_of_positive_physical_sector_factorization}.
\end{proof}
